\documentclass[11pt]{article}
\usepackage[margin=1in]{geometry}
\usepackage{amsmath,amssymb,amsthm,hyperref}
\newtheorem{theorem}{Theorem}
\newtheorem{corollary}[theorem]{Corollary}
\title{Characteristic polynomials of multiweighted partitions}
\author{Wentao Li}
\date{November 20, 2026}
\begin{document}
\maketitle
\begin{abstract}
Gonz\'alez D'Le\'on and Wachs ask for a formula for the equal-variable specialization
of the weighted bond symmetric function of a complete graph. We express it by coefficient
extraction from a power of a Laguerre polynomial. The reversed polynomial is a moment
of a translated sum of independent gamma variables. For three or more colors its number
of real zeros is zero or one according to the parity of the ground-set size. Its absolute
coefficients are strictly log-concave.
\end{abstract}

\section{The two coefficient orders}
Fix $n\ge1$ and $r\ge2$. In the multiweighted partition poset on $[n]$ a block $C$
carries a vector in $\mathbb N_0^r$ of sum $|C|-1$. Moving up the poset merges
blocks. In each resulting block the new vector must dominate the sum of its
constituents' vectors coordinatewise. The rank of a partition $\pi$ with $k$ blocks
is $n-k$. Write $\mu(\pi)=\mu(\hat 0,\pi)$ and put
\[
 F_{n,r}(t)=\sum_\pi\mu(\pi)t^{n-|\pi|},\qquad
 P_{n,r}(x)=\sum_\pi\mu(\pi)x^{|\pi|-1}.
\]
Question 7.1 of \cite{bond} asks for the equal-variable specialization when $r>2$.
Equation (7.1) of that paper gives $F$ after specialization.
Its displayed examples (7.3) and (7.4) use the coefficient order of $P$.
The two conventions are related by $F(t)=t^{n-1}P(1/t)$.

\begin{theorem}\label{main}
Let $Q_r(u)=\sum_{j=0}^{r-2}\binom{r-2}{j}u^j/(j+1)!$. Then
\begin{align*}
 P_{n,r}(x)&=(n-1)![u^{n-1}]\frac{\exp((x-n)u)}{Q_r(u)^n},\\
 F_{n,r}(t)&=(n-1)![u^{n-1}]\frac{\exp((1-nt)u)}{Q_r(tu)^n}.
\end{align*}
\end{theorem}
\begin{proof}
Let $b_m$ sum the M\"obius values of the one-block elements on $m$ labels and put
\[
 B(z)=\sum_{m\ge1}b_mz^m/m!.
\]
The bottom interval of a weighted partition is the
product of the bottom intervals on its blocks. Thus the sum of its M\"obius values
over all choices of weights for a fixed underlying partition is $\prod_C b_{|C|}$.

A partition with $k$ weighted blocks lies below exactly
$c_k=\binom{k+r-2}{k-1}$ one-block elements. Indeed the surplus in the top weight
is an arbitrary weak composition of $k-1$ into $r$ coordinates. For each top
element $\tau$, sum $\mu(\pi)$ over $\pi\le\tau$. This sum is zero unless
$\tau=\hat0$, which occurs precisely for $m=1$. Summing over all top elements
on $m$ labels, then summing over label sizes, yields
\[
 A(B(z))=z,\qquad A(u)=\sum_{k\ge1}c_k\,u^k/k!.
\]
Here $m![z^m]B(z)^k/k!$ is the sum of the products of $b_{|C|}$ over unordered
partitions into $k$ blocks. This is the equal-color instance of the inverse series
in \cite{multiple}, Theorem 3.2. Proposition 6.1 there gives the weighted rank sums. Marking each block by $x$ in the same enumeration yields
$xP_{n,r}(x)=n![z^n]\exp(xB(z))$. Lagrange inversion therefore gives
\[
 P_{n,r}(x)=(n-1)![u^{n-1}]e^{xu}\left(\frac{u}{A(u)}\right)^n .
\]
Vandermonde's identity gives
$k![u^k]ue^uQ_r(u)=\sum_j\binom{k}{j+1}\binom{r-2}{j}=c_k$ for every $k\ge1$.
Substituting $A(u)=ue^uQ_r(u)$ proves the first formula, and coefficient reversal
gives the second.
\end{proof}

For example, $Q_2=1$ recovers $P_{n,2}=(x-n)^{n-1}$ from \cite{two}.
For $r=3$ we get $Q_3=1+u/2$, in particular
$P_{3,3}(x)=x^2-9x+21$. The next result shows why the factorization into
real linear factors in the two-color case breaks down.

\begin{corollary}
For $r\ge3$, odd $n$ gives no real zero. Even $n$ gives exactly one real zero,
which is simple and belongs to $(n,\infty)$ for $P$ and to $(0,1/n)$ for $F$.
\end{corollary}
\begin{proof}
Put $d=r-2$. The Laguerre expansion gives $Q_r(u)=L_d^{(1)}(-u)/(d+1)$.
Orthogonality for the weight $y e^{-y}$ on $(0,\infty)$ implies that
$L_d^{(1)}$ has $d$ simple positive roots $\lambda_1,\ldots,\lambda_d$
\cite{dlmf}. Consequently $Q_r(u)=\prod_j(1+u/\lambda_j)$. If $S$ is the
sum of independent gamma variables of shape $n$ and respective rates $\lambda_j$,
then $\mathbb E e^{-uS}=Q_r(u)^{-n}$ near zero. Extracting the Taylor coefficient
in Theorem~\ref{main} gives
\[
 P_{n,r}(x)=\mathbb E (x-n-S)^{n-1}.
\]
The law of $S$ has a density on $(0,\infty)$. Positive even moments give the
odd-$n$ assertion (including the constant for $n=1$). For even $n$ the polynomial
has positive derivative everywhere, is negative at $n$, and tends to infinity
as $x\to\infty$. Its constant term is nonzero, so reversal gives the asserted
root count for $F$.
\end{proof}

\begin{corollary}
The absolute coefficients of $F_{n,r}$ are positive and strictly log-concave
at every interior index, for all $n\ge1$ and $r\ge2$.
\end{corollary}
\begin{proof}
For $0\le k\le n-1$, let $a_k=|[t^k]F_{n,r}(t)|$. The first formula of
Theorem~\ref{main} gives
\[
 a_k=\frac{(n-1)!}{(n-1-k)!}h_k,\qquad
 \sum_{k\ge0} h_k u^k=e^{nu}\prod_{j=1}^{r-2}(1-u/\lambda_j)^{-n}.
\]
Indeed the coefficient before taking absolute values is $(-1)^ka_k$.
The exponential series has positive log-concave coefficients and each geometric
series $(1-u/\lambda_j)^{-1}$ does as well. Convolution preserves log-concavity for
positive sequences supported on the nonnegative integers. For completeness, the
lower triangular Toeplitz matrix of such a log-concave sequence has all its
two-row, two-column minors nonnegative, by the decreasing ratios of adjacent terms
(with zeros above the diagonal). Products of these matrices correspond to convolution
and preserve this minor condition by Cauchy--Binet on finite truncations.
Thus $(h_k)$ is log-concave. The falling factorial prefactor has strictly decreasing
successive ratios $n-1-k$ from term $k$ to term $k+1$, proving strictness.
\end{proof}

This proves the ordinary equal-variable consequence of Conjecture 7.3 of \cite{bond}
for complete graphs. The conjecture itself asks for Schur positivity for every graph.

\paragraph{Status and credit.} Research draft prepared with assistance from OpenAI
Codex. The inverse-series identity and the two-color formula are due to the original
authors. Independent review of the proposed equal-color formula, its consequences and
its priority is pending.
\begin{thebibliography}{9}
\bibitem{bond} R.~S. Gonz\'alez D'Le\'on and M.~L. Wachs,
\emph{Weighted bond posets and a new chromatic symmetric function}, 2026, v2,
Question 7.1 and Section 7.2. \url{https://arxiv.org/abs/2608.08692v2}
\bibitem{multiple} R.~S. Gonz\'alez D'Le\'on, \emph{On the free Lie algebra with
multiple brackets}, 2014 preprint. \url{https://arxiv.org/abs/1408.5415}
\bibitem{two} R.~S. Gonz\'alez D'Le\'on and M.~L. Wachs, \emph{On the (co)homology
of the poset of weighted partitions}, Theorem 2.8.
\url{https://arxiv.org/abs/1309.5527}
\bibitem{dlmf} NIST Digital Library of Mathematical Functions, Chapter 18,
Sections 18.2(vi) and 18.5(iii).
\url{https://dlmf.nist.gov/18.2#vi} and \url{https://dlmf.nist.gov/18.5}
\end{thebibliography}
\end{document}
